%% kapa: exam=1 round=29 year=2018 period=1 form=A subject=civil_law subject_ko=민법 seq=29 number=29
%% source: Q-Net 감정평가사 기출문제 (cst00309) artlSeq=5202808 file=extracted/1008/5202808/2018년 제29회 감정평가사 1차시험 기출문제/2018년 제29회 감정평가사 1차시험 기출문제/1교시 A형.pdf
%% conversion: kordoc 4.24.51 -> tools/gen1.py (auto); review=proofread-sonnet; flags=none
\begin{question}{29}
법률행위에 의하지 않은 물권변동에 관한 설명으로 옳지 않은 것은? (다툼이 있으면 판례에 따름)
\begin{choices}
  \item 법정저당권은 저당권설정등기 없이 성립한다.
  \item 부동산소유권을 확인하는 판결에 의해서도 등기 없이 그 부동산의 소유권을 취득한다.
  \item 공경매에 있어서 부동산 물권변동의 시기는 매각허가결정이 확정된 후 매수인이 매각대금을 완납한 때이다.
  \item 자기의 비용과 노력으로 건물을 신축한 건축주는 건물의 소유권을 등기없이 취득한다.
  \item 상속에 의한 물권변동은 피상속인의 사망 시에 발생한다.
\end{choices}
\end{question}
\answer{2}
